%@ pages 61-67
%@ status translated
% Original section 6, "Le point de vue de la localisation", pp. 61--67.
% The authors' v3 remarks on the error in v1/v2 are part of the text and are
% kept in place, as are the %%% markers delimiting the passages revised in v3.
% The French prints the abbreviation "n-précats" (for "n-précatégories") and
% writes "$q: \Xx /Y\rightarrow X$" and "au-dessus de $X$" in the
% counterexample; these are rendered as printed, without emendation.

\clearpage
\section{The localisation viewpoint}
\label{localisationpage}

%%%%%%%%%%% most of this section is changed in v3 and we don't
%%%%%%%%%%% explicitly comment many of the changes %%%%%%%%%%%%%%%%%%%%%%

The passage from the closed model category structure for the coarse topology
to that for the topology $\Gg$ is a \emph{Bousfield localisation}, an operation
which is very well explained in Hirschhorn \cite{Hirschhorn}; see also
Goerss-Jardine \cite{JardineGoerssLoc}, and which of course goes back to
Bousfield-Kan \cite{BousfieldKan}. The operation of passing from the coarse
topology to the topology $\Gg$ preserves the notion of cofibration (both in the
setting of \ref{cmf} and in the setting of \ref{cmfHirschho}), so it is a left
Bousfield localisation in the terminology of \cite{Hirschhorn}.
%%%%%%%%%%%%%%%%%%%%%%%%%%%%%%%%%%%%%%%%%%%%%%%%v3ah
There was an error in our proof (of {\tt v1, v2}) of Lemma \ref{pourloc}
below. We thank D. Dugger for finding this mistake, and B. Toen for
communicating it to us. Lemma \ref{pourloc}, and therefore
Proposition \ref{naqualocaliser}, are not true in the setting of Segal
$n$-prestacks (and not even for simplicial sets, i.e. Segal $0$-prestacks). The
problem comes from the interaction between homotopy in arbitrarily large degree
and cohomology in arbitrarily large degree. As Dugger shows in his recent paper
\cite{Dugger}, one obtains correct statements by replacing ``covering sieve''
everywhere by ``hypercovering''. We do not wish to enter into the details here
and refer to \cite{Dugger} for this. We therefore restrict ourselves, for this
version {\tt v3}, to stating Proposition \ref{naqualocaliser} and
Lemma \ref{pourloc} in the setting of non-Segal $n$-prestacks (in other words,
in the ``$n$-truncated'' case, in which there is homotopy only in degree
$\leq n$).

This modification propagates to the following results:
\newline
---the proof of Corollary \ref{Gfibrestrict} has to be redone;
\newline
---the statements of the criteria \ref{newcritere}, \ref{critereaveclim}
and \ref{critere7} have to be modified;
\ref{newcritere} and the parts (c) of these last two now apply only to the
$n$-truncated case;
\newline
---the characterisation of the essential image in \ref{correlation} has to be
modified;
\newline
---a hypothesis (5) has to be added in the statement of Theorem
\ref{dansleschamps};
\newline
---the proof of \ref{descente} is slightly modified but the result remains the
same;
\newline
---and finally, the descent result for complexes \ref{cestunchamp} is from now
on stated only for complexes bounded below.

To return to the results \ref{naqualocaliser} and \ref{pourloc} of this version
{\tt v3} we therefore restrict ourselves to considering the cmc of non-Segal
$n$-precategories or $n$-prestacks. Note that one defines in a way analogous to
Theorem \ref{cmf} (resp. \ref{cmfHirschho}) a cmc of (non-Segal) $n$-prestacks;
for the case of \ref{cmf} we shall denote it $nPCh$.

If $p:\Yy \rightarrow \Xx$ is a functor and $U$ an $n$-precategory, then one
denotes by $U_{\Yy}$ the $n$-prestack $p_!(\underline{U})$, where
$\underline{U}$ is the constant $n$-prestack on $\Yy$ with value $U$.

We shall consider cofibrations of the following type: let $X\in \Xx$ and let
$\Bb \subset \Xx /X$ be a sieve; and let $U\hookrightarrow U'$ be a cofibration
of $n$-precategories. One denotes by
$$
cof(U\hookrightarrow U'; \Bb \subset
\Xx /X)
$$
the cofibration
$$
U_{\Xx /X} \cup ^{U_{\Bb}}U'_{\Bb} \rightarrow U'_{\Xx /X}.
$$


%%%%% holds only for n-truncated case in v3 %%%
\begin{proposition}
\label{naqualocaliser}
The
closed model category $nPCh$ of \ref{cmf} (resp. of \ref{cmfHirschho}) for
the topology $\Gg$, but for non-Segal $n$-prestacks,
is the \emph{left Bousfield localisation}
(\cite{Hirschhorn} Definition 3.3.1) of the closed model category
for
the coarse topology, with respect to the morphisms
$cof(U\hookrightarrow U'; \Bb \subset \Xx /X)$, with $\Bb$ a
sieve covering $X$ for $\Gg$.
\end{proposition}
{\em Proof:} One treats the structure of \ref{cmf}, the other case being
analogous.
We shall just show the key step, which is the following lemma.
This lemma
will imply the proposition, by Hirschhorn's theory \cite{Hirschhorn} \S 3.
\eop
%%% in v3 this holds only in the n-truncated case
%%%%
\begin{lemma}
\label{pourloc}
If $E\hookrightarrow E'$ is a $\Gg$-trivial cofibration of
(non-Segal) $n$-prestacks,
then there
exists a
sequence (possibly transfinite; see the discussion of
\cite{Hirschhorn}
\cite{DHK}
\cite{HoveyBook}) of cofibrations
$$
E=E_{-1} \hookrightarrow E_0
\hookrightarrow \ldots \hookrightarrow E_k \hookrightarrow \ldots
$$
with colimit denoted $E_{\infty}$, and a trivial cofibration for
the coarse topology $E_{\infty} \hookrightarrow E''$ such that
$E\rightarrow E'$ is a retraction of $E\rightarrow E''$; and such that the
cofibrations $E_i \hookrightarrow E_{i+1}$ are either coproducts of
cofibrations of the form $cof(U\hookrightarrow U'; \Bb \subset \Xx /X)$,
or
trivial cofibrations for the coarse topology.
\end{lemma}
{\em Proof:}
We shall call an \emph{adequate sequence} any sequence of cofibrations as in
the statement of the lemma, and the \emph{composite} of this sequence the
corresponding colimit
$E_{\infty}$. We shall prove that if $f:E\rightarrow E'$ is a $\Gg$-trivial
fibration then there exists an adequate sequence, together with a morphism
defined on the composite
$E_{\infty}\rightarrow E'$ which extends $f$ and which is an equivalence
for the coarse topology. This will imply the lemma (an exercise left to the
reader).

One proves this result by induction on $n$. One shows
first the induction step
and then the case $n=0$. Let $n\geq 1$ and suppose therefore the result true
for $n-1$. One may without loss suppose that $E$ and $E'$ are fibrant for the
coarse topology. For every $X\in \Xx$ and every
pair of objects $x,y\in E(X)$,
the morphism $E_{1/} (x,y)\rightarrow E'_{1/}(x,y)$ is a $\Gg$-trivial
fibration of $(n-1)$-prestacks above $\Xx /X$; it is therefore
equivalent (for the coarse topology) to the composite of an adequate
sequence. One remarks that for every coproduct of $E_{1/}(x,y)$ with a
cofibration of $(n-1)$-prestacks $V\rightarrow V'$ above $\Xx
/X$, one can define a corresponding coproduct of $E$ with $p_!\Upsilon
(V)\rightarrow p_!\Upsilon (V')$, where $p: \Xx /X \rightarrow \Xx$
(here we use the notation $\Upsilon$ introduced in \cite{limits}).

On the other hand, if $\Bb \subset \Xx /Y$ is a sieve covering $Y\in \Xx /X$
and $U\hookrightarrow U'$ a cofibration of $(n-1)$-precategories, then
$$
p_! \Upsilon (cof(U\hookrightarrow U'; \Bb \subset (\Xx /X) / Y))
$$
is identical to
$$
cof(\Upsilon (U)\hookrightarrow \Upsilon (U'); \Bb \subset \Xx /Y);
$$
so, by applying $p_!\Upsilon $ to an adequate sequence
of $(n-1)$-prestacks on $\Xx /X$, one obtains an adequate sequence
of
$n$-prestacks on $\Xx$.

Therefore, starting from the sequence of
cofibrations for $E_{1/} (x,y)\rightarrow E'_{1/}(x,y)$ one will obtain an
adequate sequence of cofibrations for
$$
E\cup ^{p_! \Upsilon (E_{1/} (x,y))}p_! \Upsilon (E'_{1/} (x,y)).
$$
One can then add
the coproducts with a cofibration $\ast _{\Bb} \rightarrow \ast
_{\Xx}$ (which has the right form
$cof (\emptyset \hookrightarrow \ast , \Bb \subset \Xx /X )$)
for each morphism $u:\ast _{\Bb} \rightarrow E$ extending to
$u':\ast _{\Xx /X}\rightarrow E'$.
Note that the result obtained admits a morphism to $E'$ which is
(thanks to ``three for the price of two'') a $\Gg$-weak equivalence.

If one iterates this construction for
every $(X,x,y)$ and every $(X,\Bb ,u,u')$, and that a countable infinity of
times, one
obtains a factorisation
$$
E\rightarrow E'' \rightarrow E'
$$
where the first morphism is a composite of cofibrations of the
desired form; and
the second is a $\Gg$ weak equivalence which induces an
equivalence for the coarse topology on all
the $E'' _{1/}(x,y)\rightarrow
E'_{1/}(x,y)$; moreover $E''\rightarrow E'$ has the lifting property
for morphisms of the form $\ast _{\Bb} \rightarrow \ast _{\Xx /X}$. This
implies that $E'' \rightarrow E'$ is an equivalence for the coarse
topology, which gives the desired result.

It remains to treat the case $n=0$, i.e. the case of presheaves of sets.
For this, one applies essentially twice more the induction step
above. In this case one could write it explicitly in a more
elementary way. We leave this to the reader.
\eop


Rezk shows in \cite{Rezk} that
the cmc of presheaves for sheaf theory is obtained,
starting from the cmc
of presheaves without topology, by the localisation above
(i.e. Rezk treats the case of non-Segal $0$-stacks). Moreover he suggests
at the end of his letter that the same thing should be true for the
Jardine structure.

The proposition also remains true for the closed model category structure of
type HBKQ. In fact, the results of
\cite{Hirschhorn} would allow one to define a notion of \emph{$M$-stack}
and a
corresponding closed model category, for every closed model
category $M$ which is cellular and left proper (cf the terminology of
\cite{Hirschhorn}). One would take the closed model category of
diagrams $M^{\Xx ^o}$ and localise with respect to all the morphisms
$cof(U\hookrightarrow U', \Bb \subset \Xx /X)$ (which unfortunately are no
longer cofibrations for this structure) obtained by letting
$U\hookrightarrow U'$ vary in the class of
cofibrations of $M$.
%%%%%%%%%%%%%%%%%%%%%%%%%%%%%%%%%%%%%%%%%%%%%%%%%%%%
%%%%%%%%%%%%%%%%%%%%%%%%%%%%%%%%%%%%%%%%%%%%%v3ah
%%%%%%%%%%%%%%%%%%%%%%%%%%%%%%%%%%%%%%%%%%%%%%%%%%%%
One can now do the analogue of
\ref{fibrestrict}, for the
topology $\Gg$.
One would like to obtain this even for Segal $n$-prestacks,
in spite of the restriction to the $n$-truncated case for the lemma
\ref{pourloc} in this {\tt v3}.

One denotes by $\Upsilon ^{\langle k \rangle }$ the
$k$-th iterate of the operation $\Upsilon$ introduced in
\cite{limits} (but applied this time to Segal $n$-precategories);
if $A$ is an $(n-k)$-(pre)category
(of Segal) then $\Upsilon ^{\langle k\rangle } A$ is an $n$-(pre)category
(of Segal) and this extends in an obvious way to prestacks.

Let $p: \Yy \rightarrow \Xx$ be a functor, and $A\rightarrow B$ a cofibration
of Segal $n$-prestacks above $\Yy$. One then obtains a cartesian
square
$$
\begin{array}{ccc}
p_!\Upsilon A & \rightarrow & p_!\Upsilon B \\
\downarrow && \downarrow \\
\Upsilon p_!A & \rightarrow & \Upsilon p_!B
\end{array}
$$
and by induction one obtains for every $k$ the cartesian square
$$
\begin{array}{ccc}
p_!\Upsilon ^{\langle k \rangle }
A & \rightarrow & p_!\Upsilon ^{\langle k \rangle }B \\
\downarrow && \downarrow \\
\Upsilon ^{\langle k \rangle }p_!A & \rightarrow
& \Upsilon ^{\langle k \rangle } p_!B
\end{array} .
$$

In particular, every coproduct with a cofibration of the form
$$
\Upsilon ^{\langle k \rangle }p_!A  \rightarrow
\Upsilon ^{\langle k \rangle } p_!B
$$
can also be interpreted as a coproduct with a cofibration
of the form
$$
p_!\Upsilon ^{\langle k \rangle}
A \rightarrow  p_!\Upsilon ^{\langle k \rangle}B .
$$

%%%%% the following holds in general but the proof had
%%%%% to be modified in v3
%%%%%
\begin{corollary}
\label{Gfibrestrict}
Let $\Xx$ be a site and $X\in \Xx$ with $p: \Xx /X\rightarrow \Xx$ the forgetful
functor. Let $\Gg$ be the topology on $\Xx$ which induces a topology (also
denoted $\Gg$) on $\Xx /X$. Then $p_!$ takes the $\Gg$-trivial cofibrations
of Segal $n$-prestacks
on $\Xx$ to $\Gg$-trivial cofibrations on $\Xx $.
In other words, $p_!$ is a left Quillen functor and $p^{\ast}$ a
right Quillen functor with respect to the structures of \ref{SeCmf} for the
topology $\Gg$. Thus $p^{\ast}$ preserves fibrant objects and
fibrations.
\end{corollary}
{\em Proof:}
One first treats the case $n=0$. Let $f:A\rightarrow B$ be an Illusie weak
equivalence of simplicial presheaves on $\Xx /X$. We claim that $p_!f$ is
an Illusie weak equivalence on $\Xx$. Let $Y\in \Xx$ and
$a\in (p_!A)Y$. Recall that this means that one has a morphism
$g: Y\rightarrow X$ and $a\in A(Y,g)$. Let (for $i\geq 1$)
$$
\eta \in \pi _i((p_!B )Y, fa) = \pi _i(B(Y,g), fa).
$$
By Illusie's property of $f$ on $\Xx / X$,
there exists an open covering of $(Y,g)$ by $(Y_j,g_j)$ such that
the restrictions of $\eta$ to the $Y_j$ come from elements of
$\pi _i (A(Y_j,g_j), fa|_{Y_j})$. This gives the same thing for when one
considers Illusie's property on $\Xx$. In the same way
if
$$
\phi ,\phi  '\in \pi _i((p_!A)Y, a) = \pi _i(A(Y,g), a)
$$
and if $f\phi = f\phi '$ in
$\pi _i((p_!B )Y, fa) = \pi _i(B(Y,g), fa)$ then
one has a covering of $(Y,g)$ by $(Y_j,g_j)$ such that
$f\phi |_{Y_j} = f\phi '|_{Y_j}$.  Moreover this covering for $\Xx / X$ is
also a covering of $Y$ in $\Xx$ and one obtains this property
for $f$ on $\Xx$. It only remains to treat the case of $\pi _0$. Let
$$
\eta \in \pi _0((p_!B)Y) = \coprod _{Y\rightarrow X}\pi _0(B Y).
$$
There exists therefore $g: Y\rightarrow X$ such that $\eta $ comes from an
element
$u$ of $\pi _0(B(Y,g))$. There exists therefore a covering of $(Y,g)$ such that
the
$u|_{Y_j}$ come from $A(Y_j,g_j)$ and in particular $\eta|_{Y_j}$ comes
from $p_!A(Y_j)$ since it comes from the component $A(Y_j,g_j)$.
If
$$
\phi , \phi '\in \pi _0((p_!A)Y) = \coprod _{Y\rightarrow X}\pi _0(A Y)
$$
and $f\phi =f\phi '$ then in particular $f\phi$ and $f\phi '$ are in the same
component $\pi _0(B Y)$ of $(p_!B)Y$, which implies that $\phi$ and $\phi'$
are in the same component $\pi _0(AY)$ of $\pi _0(p_!A)(Y)$. The same argument
as above still allows one to conclude.

Now note that the procedure of the proof of \ref{pourloc} allows, in the
case of Segal $n$-prestacks, to decompose a morphism $E\rightarrow E'$
which is a $\Gg$-weak equivalence into a sequence of coproducts with
cofibrations of the form
$$
q_!\Upsilon ^{\langle k \rangle } (F\hookrightarrow F')
$$
with $q: \Xx /Y\rightarrow X$ and
where the $F\hookrightarrow F'$ are either of the form
$$
F = \ast _{\Bb} \hookrightarrow \ast _{\Xx /Y} = F'
$$
where $\Bb$ is a sieve covering $Y$, or
trivial Illusie cofibrations
of simplicial presheaves in the case $k=n$. For this
one uses the principle of
commutation of $q_!$ with $\Upsilon$ for coproducts established before the
present statement.
On the other hand, every cofibration of this form
$q_!\Upsilon ^{\langle k \rangle } (F\hookrightarrow F')$ is a
$\Gg$-equivalence on $\Xx$.

If $f:E\rightarrow E'$ is a $\Gg$-equivalence above $\Xx / X$ then
let us decompose it as in the preceding paragraph (on the site $\Xx /X$).
If $p: \Xx /X\rightarrow \Xx$ is the projection then for $q: Y\rightarrow X$
one obtains $p_! q_! = (pq)_!$. The image of the decomposition of
$f$ under $p_!$ is therefore a decomposition of the same form for $p_!f$.
A morphism which decomposes in this way is a $\Gg$-equivalence, so
$p_!f$ is a $\Gg$-equivalence on $\Xx$.

This proves that $p_!$ is a left
Quillen functor. It follows that its adjoint $p^{\ast}$ is a right
Quillen functor, i.e. preserves fibrations.
\eop

{\em Remark:} This corollary is valid also for the HBKQ structure.

In the HBKQ structure one also has a converse
(we think that this result
remains true for the case of Segal $n$-prestacks, but ({\tt v3})
we state it
only for the $n$-truncated case):

\begin{corollary}
\label{restrpourhbkq}
A morphism $f: A\rightarrow B$ of non-Segal $n$-prestacks
is $\Gg$-fibrant for the HBKQ structure
if and only if $A|_{\Xx /X} \rightarrow B|_{\Xx /X}$ is $\Gg$-fibrant for
HBKQ for every $X\in \Xx$.
\end{corollary}
{\em Proof:}
One direction is provided by the preceding corollary (for the
HBKQ structure). For the other direction, suppose that $f|_{\Xx /X}$ is a
$\Gg$-fibration of HBKQ for all $X$. Then $f$ has the lifting
property for the cofibrations which are trivial object by object
(in fact, according to Hirschhorn's construction
\cite{Hirschhorn}, the trivial cofibrations
for the coarse topology are retracts of composites of sequences
of cofibrations obtained by free addition of cells corresponding
to
trivial cofibrations in the base cmc $M$); and $f$ satisfies the
lifting property for the cofibrations
$cof_{\Xx }(U\hookrightarrow U'; \Bb \subset \Xx /Y)$. By the analogue of
Lemma \ref{pourloc} for the HBKQ structure, one obtains that $f$ is a
$\Gg$-fibration of HBKQ.
\eop

{\bf Counterexample:} Corollary \ref{restrpourhbkq} is not true for the
structure of \ref{cmf} (for the Segal case at least),
if the category $\Xx$ has no final object; and this
even for the coarse topology. We give a counterexample with $n=0$,
i.e. for simplicial presheaves. The underlying category $\Xx$
will be the
category with a single object $x$, and a group $\zz$ of automorphisms of
$x$. A simplicial presheaf over $X$ is just a simplicial
set with an action of $\zz$, i.e. a pair $(B,t)$ where $B$ is a simplicial
set and $t$ an automorphism of $B$. Suppose that $t$ acts without
fixed point at each level $B_k$. If $(A,s)$ is another object, and if
the automorphism $s$ has a fixed point in some $A_k$, then there exists no
morphism $(A,s)\rightarrow (B,t)$. This leads to a contradiction with
the property of Corollary \ref{restrpourhbkq}. Indeed,
the category $\Xx /x$ is the disjoint union of copies of the
trivial category,
so $(B,t)|_{\Xx /x}$ is fibrant if and only if $B$ is a Kan-fibrant
simplicial set. Now there exists such a $B$
with $t$ acting without
fixed point---for example one takes the total singular complex $Sing(\rr )$
with the translation which acts on $\rr$ by $t: b\mapsto b+1$. Now one can
manufacture a trivial cofibration
$$
g:(B,t)\hookrightarrow (A,s)
$$
such that $(A,s)$ has fixed points, to complete the example. One takes
$A=Sing (\rr ^2)$ with the automorphism $s: (a,b)\mapsto (a,b+a)$ and the
morphism
$B\rightarrow A$ defined by $g:b\mapsto (1,b)$. Note that $A$ and $B$ are
contractile, so $g$ is a trivial cofibration for the structure \ref{cmf}
(the condition for being a cofibration is just being injective).
Now the fixed locus of $A$ is the subcomplex $Sing(\{ 0\} \times \rr
)\subset Sing
(\rr ^2)$. In particular there exists no map $(A,s)\rightarrow (B,t)$
which induces the identity on $(B,t)$. This shows that $(B,t)$ is not fibrant
for the structure of \ref{cmf}.

\subnumero{Boolean localisation}

Let us mention the existence of Jardine's paper on
``Boolean localisation'' \cite{JardineBool}: it is probable that one could
give a treatment of the cmc of \ref{cmf} with the technique of
\cite{JardineBool} but we have not explored this route.

%%%%%%%%%%%%%%%%%%%%%%%v3ah
%%%%%%%%%%%%%%%%%%%%
